A set $A\subset\mathbb N$ (here $\mathbb N=\{1,2,3,\dots\}$) has the \emph{translation property} if for every integer $n\ge 1$ there exists an integer $t_n\ge 1$ such that for all integers $a$ with $1\le a\le n$,
\[
 a\in A\ \Longleftrightarrow\ a+t_n\in A.
\]
Equivalently, the membership pattern of $A$ on $\{1,2,\dots,n\}$ repeats on the translated block $\{1+t_n,2+t_n,\dots,n+t_n\}$.

Questions posed:

(Q1) Does the set $A_{2\square}:=\{m\in\mathbb N: m=x^2+y^2\text{ for some integers }x,y\}$ have the translation property?

(Q2) If the primes are partitioned into $P\sqcup Q$ with each part containing $\gg x/\log x$ primes $\le x$ for all large $x$, does the set
\[
A_P:=\{m\in\mathbb N: \text{every prime divisor of }m\text{ lies in }P\}
\]
have the translation property?

(Q3) Let $A_{\mathrm{sqf}}$ be the set of squarefree integers. Let $t_n^*$ be the minimal positive integer $t$ witnessing the translation property for $n$, i.e.
\[
 t_n^*:=\min\{t\ge 1: \forall a\le n,\ a\in A_{\mathrm{sqf}}\iff a+t\in A_{\mathrm{sqf}}\}.
\]
How fast can $t_n^*$ grow? In particular, is it true that $t_n^*>\exp(n^c)$ for some constant $c>0$?
